% =====================================================================
% Corda whitepaper finish — the house style for /academic-format.
%
% Classic R3 / Corda look: article class, a4paper, ~26mm margins,
% Latin Modern (Computer Modern) serif — the pandoc/xelatex default, so
% no mainfont is set; that already matches the Corda body font. Section
% headings stay black; the table of contents, cross-references, citations
% and URLs render in Corda navy. Title and abstract sit alone on page 1,
% the Contents gets its own page, then the body begins on a fresh page.
%
% Pass to pandoc via --include-in-header. Self-contained: it loads every
% package it depends on, and each load is idempotent against the packages
% pandoc's default LaTeX template already pulls in.
% =====================================================================

\usepackage{etoolbox}
\usepackage{xcolor}
\usepackage{booktabs}
\usepackage{microtype}

% Corda navy. The export utility points pandoc's colorlinks variables at this
% colour (linkcolor/citecolor/urlcolor/toccolor=cordablue), so the TOC,
% cross-references, citations and URLs render navy while body text and section
% headings stay black — as in the Corda paper. Defining it here (the header is
% emitted before pandoc's own \hypersetup) is what lets those variables resolve.
\definecolor{cordablue}{HTML}{1A2A6C}

% Isolate the title-and-abstract page, then give the Contents its own page.
\preto\tableofcontents{\clearpage}
\appto\tableofcontents{\clearpage}

% Corda tables breathe: a touch more row height, no vertical rules
% (booktabs convention, reinforced here for any hand-written tabulars).
\renewcommand{\arraystretch}{1.15}
